\documentclass{article}
\usepackage{iftex}
\ifXeTeX
  \usepackage{fontspec}
  \usepackage{xeCJK}
  \usepackage[haranoaji]{zxjafont}
  \XeTeXlinebreaklocale ""
\fi
\ifLuaTeX
  \usepackage[haranoaji]{luatexja-preset}
\fi
\usepackage[a4paper, margin=2cm]{geometry}
\usepackage{graphicx}
\newdimen\charwd
{\tt\global\setbox0=\hbox{x}\global\charwd=\wd0}
\frenchspacing
\usepackage{fancyhdr}
\pagestyle{fancy}
\renewcommand{\headrulewidth}{0pt}
\fancyhf{}
\fancyhead[R]{\rm src2tex-go version 2.23}
% plain TeX compatibility
\makeatletter
\providecommand{\eqalign}[1]{\vcenter{\openup1\jot\m@th
  \ialign{\strut\hfil$\displaystyle{##}$&$\displaystyle{{}##}$\hfil\crcr#1\crcr}}}
\makeatother
\fancyfoot[R]{\rm\hfill samples{\tt /}newton.c\qquad page \thepage}
\begin{document}

\ifx\sevenrm\undefined
  \font\sevenrm=cmr7 scaled \magstep0
\fi

\newread\MyStyle
\openin\MyStyle=src2latex.s2t
\ifeof\MyStyle
  \closein\MyStyle
\else
  \input src2latex.s2t
  \closein\MyStyle
\fi

\def\mc{\relax}
\def\gt{\relax}
\def\sc{\scshape}

\tt\mc 

\noindent
\rm\mc {\tt /}{\tt *}\kern1\charwd  \hrulefill \rm\mc \kern1\charwd {\tt *}

\noindent
\mbox{}\hfill

\noindent
\mbox{} \bf Newton-Raphson 法 \rm\mc 

\noindent
\mbox{}\hfill

\noindent
\mbox{}方程式\kern1\charwd $f(x)=0$\rm\mc \kern1\charwd を解くためには、漸化式

\noindent
\mbox{}\hfill

\noindent
\mbox{}$\displaystyle\qquad
	x_0=a,\quad x_{n+1}=x_n-{f(x_n)\over f'(x_n)}
	\qquad\cdots\cdots\ (\star)$\rm\mc 

\noindent
\mbox{}\hfill

\noindent
\mbox{}を、適当な初期値\kern1\charwd $a$\rm\mc \kern1\charwd に対して解いて、数列\kern1\charwd $\{x_n\}$\rm\mc \kern1\charwd を構成する。

\noindent
\mbox{}初期値\kern1\charwd $a$\rm\mc \kern1\charwd がうまく与えられれば、この数列は上述の方程式の

\noindent
\mbox{}１つの解\kern1\charwd $\alpha$\rm\mc \kern1\charwd に、収束することが知られている,\kern1\charwd  \it i.e. \rm\mc ,

\noindent
\mbox{}\hfill

\noindent
\mbox{}$\displaystyle\qquad
	\exists\ \alpha=\lim_{n\to\infty}x_n
	\quad such\ that\quad f(\alpha)=0\ .$\rm\mc 

\noindent
\mbox{}\hfill

\noindent
\mbox{}漸化式\kern1\charwd $(\star)$\rm\mc \kern1\charwd の意味とその収束の様子は、次の図と式を見れば

\noindent
\mbox{}一目瞭然である。

\noindent
\mbox{}\hfill

\noindent
\begin{center}
\includegraphics[scale=.7]{newton}
\end{center}
 \rm\mc 

\noindent
\mbox{}\hfill

\noindent
\mbox{}$\displaystyle\qquad
	y-f(x_n)=f'(x_n)(x-x_n)
	\qquad\cdots\cdots\ \hbox{点}\ (x_n,f(x_n))\ \hbox{を通る接線の方程式}$\rm\mc 

\noindent
\mbox{}\hfill

\noindent
\mbox{}$\displaystyle\qquad
	x_n-{f(x_n)\over f'(x_n)}
	\qquad\cdots\cdots\ \hbox{上記の接線と}\ x\ \hbox{軸との交点の}
	\ x\ \hbox{座標}$\rm\mc 

\noindent
\mbox{}\hfill

\noindent
\mbox{}数列\kern1\charwd $\{x_n\}$\rm\mc \kern1\charwd の収束性の証明および収束のオーダーの評価は、

\noindent
\mbox{}それほど簡単なことではない。$\epsilon$\rm\mc {\tt -}$\delta$\rm\mc \kern1\charwd 論法と解析学に関する

\noindent
\mbox{}知識が必要とされる。興味のある学生には、参考文献を紹介する。

\noindent
\mbox{}\hfill

\noindent
\mbox{}以下のCソース\kern1\charwd newton.c\kern1\charwd は、方程式

\noindent
\mbox{}\hfill

\noindent
\mbox{}$\qquad\displaystyle x^2-5=0$\rm\mc 

\noindent
\mbox{}\hfill

\noindent
\mbox{}の解の１つが\kern1\charwd $\sqrt{5}$\rm\mc \kern1\charwd であることに注目して、Newton{\tt -}Raphson\kern1\charwd 法で\kern1\charwd $\sqrt{5}$\rm\mc 

\noindent
\mbox{}を求めるプログラムである。 \tt A, F(X), DF(X) \rm\mc \kern1\charwd の定義をいろいろ

\noindent
\mbox{}と変えて、各人で数値実験を行ってみよう。

\noindent
\mbox{}\hfill

\noindent
\mbox{}{\tt *}\kern1\charwd  \hrulefill \rm\mc \kern1\charwd {\tt *}{\tt /}
\tt\mc 

\noindent
\mbox{}\hfill

\noindent
\mbox{}\rm\mc {\tt /}{\tt *}\kern1\charwd  \hrulefill\ newton.c\ \hrulefill \rm\mc \kern1\charwd {\tt *}{\tt /}
\tt\mc 

\noindent
\mbox{}\hfill

\noindent
\mbox{}\hfill

\noindent
\mbox{}{\tt\#}include{\tt\mc \kern1\charwd}{\tt <}stdio.h{\tt >}

\noindent
\mbox{}{\tt\#}define{\tt\mc \kern1\charwd}A{\tt\mc \kern1\charwd}4.{\tt\mc \kern4\charwd}{\tt\mc \kern8\charwd}{\tt\mc \kern8\charwd}{\tt\mc \kern8\charwd}\rm\mc {\tt /}{\tt *}\kern1\charwd 初期値\kern1\charwd  \hfill \rm\mc \kern1\charwd {\tt *}{\tt /}
\tt\mc 

\noindent
\mbox{}{\tt\#}define{\tt\mc \kern1\charwd}F(X){\tt\mc \kern1\charwd}((X){\tt *}(X){\tt -}5.){\tt\mc \kern7\charwd}{\tt\mc \kern8\charwd}\rm\mc {\tt /}{\tt *}\kern1\charwd 与えられた関数\kern1\charwd  \hfill \rm\mc \kern1\charwd {\tt *}{\tt /}
\tt\mc 

\noindent
\mbox{}{\tt\#}define{\tt\mc \kern1\charwd}DF(X){\tt\mc \kern1\charwd}(2.{\tt *}(X)){\tt\mc \kern2\charwd}{\tt\mc \kern8\charwd}{\tt\mc \kern8\charwd}\rm\mc {\tt /}{\tt *}\kern1\charwd その導関数\kern1\charwd  \hfill \rm\mc \kern1\charwd {\tt *}{\tt /}
\tt\mc 

\noindent
\mbox{}\hfill

\noindent
\mbox{}main()

\noindent
\mbox{}{\tt\char'173}

\noindent
\mbox{}{\tt\mc \kern4\charwd}double{\tt\mc \kern1\charwd}x,{\tt\mc \kern1\charwd}y;{\tt\mc \kern8\charwd}{\tt\mc \kern8\charwd}{\tt\mc \kern8\charwd}\rm\mc {\tt /}{\tt *}\kern1\charwd 倍精度で計算する\kern1\charwd  \hfill \rm\mc \kern1\charwd {\tt *}{\tt /}
\tt\mc 

\noindent
\mbox{}\hfill

\noindent
\mbox{}{\tt\mc \kern4\charwd}x{\tt\mc \kern1\charwd}={\tt\mc \kern1\charwd}A;{\tt\mc \kern6\charwd}{\tt\mc \kern8\charwd}{\tt\mc \kern8\charwd}{\tt\mc \kern8\charwd}\rm\mc {\tt /}{\tt *}\kern1\charwd $\displaystyle x_0=a$\rm\mc \kern1\charwd  \hfill \rm\mc \kern1\charwd {\tt *}{\tt /}
\tt\mc 

\noindent
\mbox{}{\tt\mc \kern4\charwd}printf({\tt "}{\tt\%}.16f{\tt\char92}n{\tt "},{\tt\mc \kern1\charwd}x);{\tt\mc \kern7\charwd}{\tt\mc \kern8\charwd}\rm\mc {\tt /}{\tt *}\kern1\charwd $x_0$\rm\mc \kern1\charwd の表示\kern1\charwd  \hfill \rm\mc \kern1\charwd {\tt *}{\tt /}
\tt\mc 

\noindent
\mbox{}{\tt\mc \kern4\charwd}y{\tt\mc \kern1\charwd}={\tt\mc \kern1\charwd}x{\tt\mc \kern1\charwd}{\tt -}{\tt\mc \kern1\charwd}F(x){\tt\mc \kern1\charwd}{\tt /}{\tt\mc \kern1\charwd}DF(x);{\tt\mc \kern1\charwd}{\tt\mc \kern6\charwd}{\tt\mc \kern8\charwd}\rm\mc {\tt /}{\tt *}\kern1\charwd $\displaystyle
				   	x_1=x_0 - {f(x_0)\over f'(x_0)}$\rm\mc 

\noindent
\mbox{}{\tt\mc \kern8\charwd}{\tt\mc \kern8\charwd}{\tt\mc \kern8\charwd}{\tt\mc \kern8\charwd}{\tt\mc \kern8\charwd} \hfill \rm\mc \kern1\charwd {\tt *}{\tt /}
\tt\mc 

\noindent
\mbox{}{\tt\mc \kern4\charwd}printf({\tt "}{\tt\%}.16f{\tt\char92}n{\tt "},{\tt\mc \kern1\charwd}y);{\tt\mc \kern7\charwd}{\tt\mc \kern8\charwd}\rm\mc {\tt /}{\tt *}\kern1\charwd $x_1$\rm\mc \kern1\charwd の表示\kern1\charwd  \hfill \rm\mc \kern1\charwd {\tt *}{\tt /}
\tt\mc 

\noindent
\mbox{}{\tt\mc \kern4\charwd}while{\tt\mc \kern1\charwd}(x{\tt\mc \kern1\charwd}!={\tt\mc \kern1\charwd}y){\tt\mc \kern6\charwd}{\tt\mc \kern8\charwd}{\tt\mc \kern8\charwd}\rm\mc {\tt /}{\tt *}\kern1\charwd  {}iteration をやっても
					   値が変わらなくなったら終了
					   \hfill \rm\mc \kern1\charwd {\tt *}{\tt /}
\tt\mc 

\noindent
\mbox{}{\tt\mc \kern4\charwd}{\tt\char'173}

\noindent
\mbox{}{\tt\mc \kern8\charwd}x{\tt\mc \kern1\charwd}={\tt\mc \kern1\charwd}y;

\noindent
\mbox{}{\tt\mc \kern8\charwd}y{\tt\mc \kern1\charwd}={\tt\mc \kern1\charwd}x{\tt\mc \kern1\charwd}{\tt -}{\tt\mc \kern1\charwd}F(x){\tt\mc \kern1\charwd}{\tt /}{\tt\mc \kern1\charwd}DF(x);{\tt\mc \kern1\charwd}{\tt\mc \kern2\charwd}{\tt\mc \kern8\charwd}\rm\mc {\tt /}{\tt *}\kern1\charwd $\displaystyle x_{n+1}=x_n
				   	- {f(x_n)\over f'(x_n)}$\rm\mc \kern1\charwd  \hfill \rm\mc \kern1\charwd {\tt *}{\tt /}
\tt\mc 

\noindent
\mbox{}{\tt\mc \kern8\charwd}printf({\tt "}{\tt\%}.16f{\tt\char92}n{\tt "},{\tt\mc \kern1\charwd}y);{\tt\mc \kern3\charwd}{\tt\mc \kern8\charwd}\rm\mc {\tt /}{\tt *}\kern1\charwd $x_{n+1}$\rm\mc \kern1\charwd の表示\kern1\charwd  \hfill \rm\mc \kern1\charwd {\tt *}{\tt /}
\tt\mc 

\noindent
\mbox{}{\tt\mc \kern4\charwd}{\tt\char'175}

\noindent
\mbox{}{\tt\char'175}

\noindent
\mbox{}

\rm\mc

\end{document}
